% TOP_PROBLEM: {"id": 30007091, "title": "Power improvement for planar halving lines", "category_id": 2, "category": {"id": 2, "name": "combinatorics", "display_name": "Combinatorics", "description": "Counting problems, graph theory, discrete structures.", "slug": "combinatorics", "order_index": 2, "created_at": "2026-07-31T15:26:25.671Z"}, "status": "open"}
% FIRST_POSTED: 2026-10-06
% ENTRY_KIND: statement_only
% !TeX program = lualatex
\documentclass[11pt]{article}
\usepackage[margin=25mm]{geometry}
\usepackage{fontspec}
\setmainfont{Latin Modern Roman}
\usepackage{amsmath,amssymb,mathtools}
\usepackage{unicode-math}
\setmathfont{Latin Modern Math}
\usepackage{xurl,hyperref}
\hypersetup{colorlinks=true,urlcolor=blue}
\setcounter{secnumdepth}{0}
\setlength{\parindent}{0pt}
\setlength{\parskip}{5pt}
\setlength{\emergencystretch}{3em}
\newcommand{\R}{\mathbb R}
\newcommand{\N}{\mathbb N}
\newcommand{\Z}{\mathbb Z}
\newcommand{\Q}{\mathbb Q}
\newcommand{\C}{\mathbb C}
\newcommand{\D}{\mathbb D}
\newcommand{\F}{\mathbb F}
\newcommand{\sref}[2]{\hyperlink{source:#1:#2}{[S#2]}}
\newcommand{\cataloguescope}[1]{#1}
\begin{document}
\section*{Power improvement for planar halving lines}\label{30007091}
\subsection{Definitions and mathematical statement}
\paragraph{Statement recorded in the supplied source.}
Do absolute constants $\epsilon >0$ and C exist such that every even n and every n-point planar set in general position determine at most C $n^{4/3-\epsilon }$ lines that leave exactly (n-2)/2 other points on each side?
\subsection{Short English statement}
Do absolute constants $\epsilon >0$ and C exist such that every even n and every n-point planar set in general position determine at most C $n^{4/3-\epsilon }$ lines that leave exactly (n-2)/2 other points on each side?
\subsection{Sources}
\begin{itemize}
\item[S1] Original question source or mathematical context.
\url{https://doi.org/10.1007/PL00009354}.
\item[S2] Original question source or mathematical context.
\url{https://doi.org/10.1007/PL00009354}.
\item[S3] Original question source or mathematical context.
\url{https://users.renyi.hu/~p_erdos/1973-07.pdf}.
\item[S4] Original question source or mathematical context.
\url{https://users.renyi.hu/~p_erdos/1973-07.pdf}.
\item[S5] OpenAI, manuscript citation page. Source of the formulation motivating this local problem.
\url{https://github.com/openai/math/blob/adc7f1241b42e322a6451854ab7e4b4c146bf78a/preprints/A-power-saving-for-planar-halving-lines-September-25-2026/README.md}.
\end{itemize}
\end{document}
